\documentclass[10pt]{article}
\usepackage[T1,T2A]{fontenc}
\usepackage[utf8]{inputenc}
\usepackage[english,ukrainian]{babel}
\usepackage{geometry}
\usepackage{longtable}
\usepackage{array}
\usepackage{multirow}
\usepackage{hyperref}
\usepackage{mathptmx}
\renewcommand{\familydefault}{\rmdefault}

\geometry{a4paper,left=2cm,right=2cm,top=2cm,bottom=2cm}

\begin{document}

\begin{flushright}
ЗАТВЕРДЖЕНО\\
Наказ Держспецзв’язок України\\
\_\_\_\_\_\_\_\_\_\_\_\_ 2026 року №\_\_\_\_
\end{flushright}

\begin{center}
\textbf{\large Базовий профіль безпеки системи, де обробляється службова інформація, затверджений наказом Адміністрації Держспецзв’язку від 02.07.2025 № 419}
\end{center}

\footnotesize
\begin{longtable}{|c|>{\raggedright\arraybackslash}p{4cm}|>{\raggedright\arraybackslash}p{6cm}|>{\raggedright\arraybackslash}p{1.5cm}|>{\raggedright\arraybackslash}p{3.5cm}|}
\hline
\textbf{№ з/п} & \textbf{Назва дії з безпеки інформації} & \textbf{Зміст дії} & \textbf{Заходи захисту} & \textbf{Мінімальні необхідні параметри}\\
\hline
\endfirsthead

\hline
\textbf{№ з/п} & \textbf{Назва дії з безпеки інформації} & \textbf{Зміст дії} & \textbf{Заходи захисту} & \textbf{Мінімальні необхідні параметри}\\
\hline
\endhead

\hline
\endfoot

\hline
\endlastfoot

\multicolumn{5}{|>{\raggedright\arraybackslash}p{16cm}|}{\textbf{УПРАВЛІННЯ ДОСТУПОМ (AC)}} \\* \noalign{\hrule height 0.4pt}
1 & УПРАВЛІННЯ ОБЛІКОВИМИ ЗАПИСАМИ (AC-2) & визначено персонал або ролі, необхідні для затвердження запитів на створення облікових записів & AC-2 & Параметр: ac\_\allowbreak{}2\_\allowbreak{}odp\_\allowbreak{}03\newline Тип: list\newline Значення: \textbf{admin, security\_\allowbreak{}officer} \\ \cline{3-5}
 &  & визначено персонал або ролі, які мають бути повідомлені &  & Параметр: ac\_\allowbreak{}2\_\allowbreak{}odp\_\allowbreak{}05\newline Тип: list\newline Значення: \textbf{admin, security\_\allowbreak{}officer} \\ \cline{3-5}
 &  & визначено період часу, протягом якого адміністратори облікових записів повинні бути повідомлені про те, що облікові записи більше не потрібні &  & Параметр: ac\_\allowbreak{}2\_\allowbreak{}odp\_\allowbreak{}06\newline Тип: integer\newline Значення: \textbf{30} \\ \cline{3-5}
 &  & визначено термін, протягом якого необхідно повідомляти адміністраторів облікових записів про звільнення або переведення користувачів &  & Параметр: ac\_\allowbreak{}2\_\allowbreak{}odp\_\allowbreak{}07\newline Тип: integer\newline Значення: \textbf{30} \\ \cline{3-5}
 &  & визначено період часу, протягом якого необхідно повідомляти адміністраторів облікових записів про зміни у використанні системи або необхідність знати про зміни для окремої особи &  & Параметр: ac\_\allowbreak{}2\_\allowbreak{}odp\_\allowbreak{}08\newline Тип: integer\newline Значення: \textbf{30} \\ \cline{3-5}
 &  & визначено періодичність перегляду облікових записів &  & Параметр: ac\_\allowbreak{}2\_\allowbreak{}odp\_\allowbreak{}10\newline Тип: integer\newline Значення: \textbf{30} \\ \cline{3-5}
 &  & Визначено період часу, протягом якого необхідно деактивувати облікові записи & AC-2(3) & Параметр: ac\_\allowbreak{}2\_\allowbreak{}3\_\allowbreak{}odp\_\allowbreak{}01\newline Тип: integer\newline Значення: \textbf{30} \\ \cline{3-5}
 &  & Визначено період часу неактивності, після закінчення якого облікові записи будуть деактивовані &  & Параметр: ac\_\allowbreak{}2\_\allowbreak{}3\_\allowbreak{}odp\_\allowbreak{}02\newline Тип: integer\newline Значення: \textbf{30} \\ \cline{3-5}
 &  & Користувачі повинні виходити з системи, коли період очікуваної бездіяльності або опис часу, коли потрібно вийти з системи & AC-2(5) & Параметр: ac\_\allowbreak{}2\_\allowbreak{}5\_\allowbreak{}01\newline Тип: integer\newline Значення: \textbf{30} \\ \cline{3-5}
 &  & Визначено часовий період очікуваної бездіяльності або опис, коли потрібно вийти з системи &  & Параметр: ac\_\allowbreak{}2\_\allowbreak{}5\_\allowbreak{}odp\newline Тип: integer\newline Значення: \textbf{30} \\ \cline{3-5}
 &  & Облікові записи користувачів деактивуються протягом періоду часу з моменту виявлення значних ризиків & AC-2(13) & Параметр: ac\_\allowbreak{}2\_\allowbreak{}13\_\allowbreak{}01\newline Тип: integer\newline Значення: \textbf{30} \\ \cline{3-5}
 &  & Визначено період часу, протягом якого необхідно деактивувати облікові записи фізичних осіб, які становлять значний ризик &  & Параметр: ac\_\allowbreak{}2\_\allowbreak{}13\_\allowbreak{}odp\_\allowbreak{}01\newline Тип: integer\newline Значення: \textbf{30} \\ \hline
2 & ЗАБЕЗПЕЧЕННЯ ДОСТУПУ (AC-3) & Затверджені повноваження на логічний доступ до інформації та ресурсів системи виконуються відповідно до чинних політик(правил) управління доступом & AC-3 & Параметр: ac\_\allowbreak{}3\_\allowbreak{}01\newline Тип: list\newline Значення: \textbf{default\_\allowbreak{}deny\_\allowbreak{}rule, abac\_\allowbreak{}rule\_\allowbreak{}1} \\ \hline
3 & УПРАВЛІННЯ ІНФОРМАЦІЙНИМИ ПОТОКАМИ (AC-4) & Затверджені повноваження застосовуються для контролю потоку інформації всередині системи та між підключеними системами на основі політики управління інформаційними потоками & AC-4 & Параметр: ac\_\allowbreak{}4\_\allowbreak{}01\newline Тип: list\newline Значення: \textbf{default\_\allowbreak{}deny\_\allowbreak{}rule, abac\_\allowbreak{}rule\_\allowbreak{}1} \\ \cline{3-5}
 &  & Визначено політики управління інформаційними потоками всередині системи та між підключеними системами &  & Параметр: ac\_\allowbreak{}4\_\allowbreak{}odp\newline Тип: list\newline Значення: \textbf{default\_\allowbreak{}deny\_\allowbreak{}rule, abac\_\allowbreak{}rule\_\allowbreak{}1} \\ \hline
4 & РОЗМЕЖУВАННЯ ОБОВ'ЯЗКІВ (AC-5) & & AC-5 & \\ \hline
5 & МІНІМІЗАЦІЯ ПОВНОВАЖЕНЬ (AC-6) & & AC-6 & \\ \cline{3-5}
 &  & & AC-6(1) & \\ \cline{3-5}
 &  & Повноваження, призначені ролям і класам, переглядаються з частотою для перевірки необхідності таких повноважень & AC-6(7) & Параметр: ac\_\allowbreak{}6\_\allowbreak{}7\_\allowbreak{}a\newline Тип: integer\newline Значення: \textbf{30} \\ \cline{3-5}
 &  & Визначено частоту перегляду повноважень, призначених ролям або класам користувачів &  & Параметр: ac\_\allowbreak{}6\_\allowbreak{}7\_\allowbreak{}odp\_\allowbreak{}01\newline Тип: integer\newline Значення: \textbf{30} \\ \cline{3-5}
 &  & Визначено ролі або класи користувачів, яким призначено повноваження  &  & Параметр: ac\_\allowbreak{}6\_\allowbreak{}7\_\allowbreak{}odp\_\allowbreak{}02\newline Тип: list\newline Значення: \textbf{admin, security\_\allowbreak{}officer} \\ \cline{3-5}
 &  & & AU-9(4) & \\ \hline
6 & НЕПРИВІЛЕЙОВАНИЙ ДОСТУП ДО НЕЗАХИЩЕНИХ ФУНКЦІЙ (AC-6(2)) (AC-6) & Користувачі облікових записів (або ролей) системи з доступом до функцій безпеки або інформації, що стосується безпеки, повинні використовувати непривілейовані облікові записи або ролі під час доступу до незахищених функцій & AC-6(2) & Параметр: ac\_\allowbreak{}6\_\allowbreak{}2\_\allowbreak{}01\newline Тип: list\newline Значення: \textbf{admin, security\_\allowbreak{}officer} \\ \cline{3-5}
 &  & Привілейовані облікові записи в системі обмежено персонал або ролі & AC-6(5) & Параметр: ac\_\allowbreak{}6\_\allowbreak{}5\_\allowbreak{}01\newline Тип: list\newline Значення: \textbf{admin, security\_\allowbreak{}officer} \\ \cline{3-5}
 &  & Визначено персонал або ролі, яким мають бути обмежені привілейовані облікові записи в системі &  & Параметр: ac\_\allowbreak{}6\_\allowbreak{}5\_\allowbreak{}odp\newline Тип: list\newline Значення: \textbf{admin, security\_\allowbreak{}officer} \\ \hline
7 & АУДИТ ВИКОРИСТАННЯ ПРИВІЛЕЙОВАНИХ ФУНКЦІЙ (AC-6(9)) (AC-6) & & AC-6(9) & \\ \cline{3-5}
 &  & & AC-6(10) & \\ \hline
8 & НЕВДАЛІ СПРОБИ ВХОДУ В СИСТЕМУ (AC-7) & Визначено кількість послідовних неуспішних спроб входу користувача, дозволених протягом певного періоду часу & AC-7 & Параметр: ac\_\allowbreak{}7\_\allowbreak{}odp\_\allowbreak{}01\newline Тип: integer\newline Значення: \textbf{3} \\ \cline{3-5}
 &  & Визначено період часу, яким обмежується кількість послідовних неуспішних спроб входу користувача &  & Параметр: ac\_\allowbreak{}7\_\allowbreak{}odp\_\allowbreak{}02\newline Тип: string\newline Значення: \textbf{15m} \\ \cline{3-5}
 &  & Період часу, на який буде заблоковано обліковий запис або вузол (якщо вибрано) &  & Параметр: ac\_\allowbreak{}7\_\allowbreak{}odp\_\allowbreak{}04\newline Тип: string\newline Значення: \textbf{15m} \\ \hline
9 & ПОПЕРЕДЖЕННЯ ПРО ВИКОРИСТАННЯ СИСТЕМИ (AC-8) & & AC-8 & \\ \hline
10 & БЛОКУВАННЯ ПРИСТРОЮ (AC-11) & Визначено часовий проміжок бездіяльності, після якого ініціюється блокування пристрою & AC-11 & Параметр: ac\_\allowbreak{}11\_\allowbreak{}odp\_\allowbreak{}02\newline Тип: string\newline Значення: \textbf{15m} \\ \cline{3-5}
 &  & & AC-11(1) & \\ \hline
11 & ПРИПИНЕННЯ СЕАНСУ (AC-12) & Визначено умови або події, що вимагають припинення сеансу & AC-12 & Параметр: ac\_\allowbreak{}12\_\allowbreak{}odp\newline Тип: list\newline Значення: \textbf{login, logout, failed\_\allowbreak{}attempt} \\ \hline
12 & ВІДДАЛЕНИЙ ДОСТУП (AC-17) & & AC-17 & \\ \cline{3-5}
 &  & & AC-17(3) & \\ \cline{3-5}
 &  & & AC-17(4) & \\ \hline
13 & БЕЗДРОТОВИЙ ДОСТУП (AC-18) & & AC-18 & \\ \cline{3-5}
 &  & Вибрано одне або декілька з наступних ЗНАЧЕНЬ ПАРАМЕТРІВ: \{користувачі; & AC-18(1) & Параметр: ac\_\allowbreak{}18\_\allowbreak{}1\_\allowbreak{}odp\newline Тип: list\newline Значення: \textbf{користувачі, пристрої} \\ 
 &  & пристрої\} &  &  \\ \cline{3-5}
 &  & Бездротовий доступ до системи захищений за допомогою шифрування &  & Параметр: ac\_\allowbreak{}18\_\allowbreak{}1\_\allowbreak{}02\newline Тип: string\newline Значення: \textbf{AES-256-GCM} \\ \cline{3-5}
 &  & & AC-18(3) & \\ \hline
14 & КОНТРОЛЬ ДОСТУПУ ДЛЯ МОБІЛЬНИХ ПРИСТРОЇВ (AC-19) & & AC-19 & \\ \cline{3-5}
 &  & Вибрано одне з наступних ЗНАЧЕНЬ ПАРАМЕТРІВ: \{повне шифрування пристроїв; & AC-19(5) & Параметр: ac\_\allowbreak{}19\_\allowbreak{}5\_\allowbreak{}odp\_\allowbreak{}01\newline Тип: list\newline Значення: \textbf{повне шифрування пристроїв, шифрування сховищ інформації} \\ 
 &  & шифрування сховищ інформації\} &  &  \\ \hline
15 & ВИКОРИСТАННЯ ЗОВНІШНІХ СИСТЕМ (AC-20) & Вибрано одне або більше з наступних ЗНАЧЕНЬ ПАРАМЕТРІВ: \{встановити <AC -20\_\allowbreak{}ODP[02] умови та положення>; & AC-20 & Параметр: ac\_\allowbreak{}20\_\allowbreak{}odp\_\allowbreak{}01\newline Тип: list\newline Значення: \textbf{умови та положення, заходи захисту} \\ 
 &  & визначити <AC-20\_\allowbreak{}ODP[03] заходи захисту>\} &  &  \\ \cline{3-5}
 &  & Визначено умови та положення, що відповідають довірчим відносинам, встановленим з іншими організаціями, які володіють, експлуатують та/або обслуговують зовнішні системи (якщо вибрано) &  & Параметр: ac\_\allowbreak{}20\_\allowbreak{}odp\_\allowbreak{}02\newline Тип: list\newline Значення: \textbf{} \\ \cline{3-5}
 &  & Авторизовані особи мають право використовувати зовнішню систему для доступу до системи або для обробки, зберігання чи передачі інформації, що контролюється організацією, лише після перевірки виконання заходів безпеки та конфіденційності, зазначених у політиці безпеки та конфіденційності організації, а також планах безпеки та конфіденційності (якщо такі застосовуються) & AC-20(1) & Параметр: ac\_\allowbreak{}20\_\allowbreak{}1\_\allowbreak{}a\newline Тип: list\newline Значення: \textbf{admin, security\_\allowbreak{}officer} \\ \cline{3-5}
 &  & Авторизовані особи мають право використовувати зовнішню систему для доступу до системи або для обробки, зберігання чи передачі інформації, що контролюється організацією, лише після збереження погоджених угод про підключення або обробку системи з структурою організації, на якій розміщена зовнішня система &  & Параметр: ac\_\allowbreak{}20\_\allowbreak{}1\_\allowbreak{}b\newline Тип: list\newline Значення: \textbf{admin, security\_\allowbreak{}officer} \\ \cline{3-5}
 &  & Використання портативних пристроїв носіїв інформації уповноваженими особами обмежено у зовнішніх системах за допомогою обмеження  & AC-20(2) & Параметр: ac\_\allowbreak{}20\_\allowbreak{}2\_\allowbreak{}01\newline Тип: list\newline Значення: \textbf{admin, security\_\allowbreak{}officer} \\ \cline{3-5}
 &  & Визначено обмеження на використання авторизованими особами портативних носіїв інформації у зовнішніх системах &  & Параметр: ac\_\allowbreak{}20\_\allowbreak{}2\_\allowbreak{}odp\newline Тип: list\newline Значення: \textbf{admin, security\_\allowbreak{}officer} \\ \hline
16 & ПУБЛІЧНО ДОСТУПНИЙ КОНТЕНТ (AC-22) & & AC-22 & \\ \hline
\multicolumn{5}{|>{\raggedright\arraybackslash}p{16cm}|}{\textbf{ОБІЗНАНІСТЬ ТА НАВЧАННЯ (AT)}} \\* \noalign{\hrule height 0.4pt}
17 & НАВЧАННЯ З ПІДВИЩЕННЯ ОБІЗНАНОСТІ (AT-2) & Визначено періодичність проведення навчання грамотності з питань безпеки для користувачів системи (в тому числі менеджерів, вищого керівництва та підрядників) після початкового тренінгу & AT-2 & Параметр: at\_\allowbreak{}2\_\allowbreak{}odp\_\allowbreak{}01\newline Тип: string\newline Значення: \textbf{щорічно} \\ \cline{3-5}
 &  & Визначено періодичність проведення навчання грамотності з питань конфіденційності для користувачів системи (в тому числі менеджерів, вищого керівництва та підрядників) після початкового тренінгу &  & Параметр: at\_\allowbreak{}2\_\allowbreak{}odp\_\allowbreak{}02\newline Тип: string\newline Значення: \textbf{щорічно} \\ \cline{3-5}
 &  & & AT-2(2) & \\ \cline{3-5}
 &  & & AT-2(3) & \\ \hline
18 & РОЛЬОВЕ НАВЧАННЯ (AT-3) & Визначено ролі та обов'язки для тренінгів з безпеки на основі ролей & AT-3 & Параметр: at\_\allowbreak{}3\_\allowbreak{}odp\_\allowbreak{}01\newline Тип: list\newline Значення: \textbf{admin, security\_\allowbreak{}officer} \\ \cline{3-5}
 &  & Визначено ролі та обов'язки для тренінгів з конфіденційності на основі ролей &  & Параметр: at\_\allowbreak{}3\_\allowbreak{}odp\_\allowbreak{}02\newline Тип: list\newline Значення: \textbf{admin, security\_\allowbreak{}officer} \\ \cline{3-5}
 &  & Визначено частоту оновлення змісту навчання на основі ролей &  & Параметр: at\_\allowbreak{}3\_\allowbreak{}odp\_\allowbreak{}04\newline Тип: integer\newline Значення: \textbf{30} \\ \hline
\multicolumn{5}{|>{\raggedright\arraybackslash}p{16cm}|}{\textbf{АУДИТ ТА ПІДЗВІТНІСТЬ (AU)}} \\* \noalign{\hrule height 0.4pt}
19 & ПОДІЇ АУДИТУ (AU-2) & Зазначені типи подій реєструються 02\_\allowbreak{}ODP[03] частота або ситуація>; & AU-2 & Параметр: au\_\allowbreak{}2\_\allowbreak{}c\_\allowbreak{}02\newline Тип: string\newline Значення: \textbf{щорічно} \\ 
 &  & <AU- &  &  \\ \cline{3-5}
 &  & Переглядаються та оновлюються типи подій, вибрані для реєстрації, частота. у системі &  & Параметр: au\_\allowbreak{}2\_\allowbreak{}e\newline Тип: string\newline Значення: \textbf{щорічно} \\ \cline{3-5}
 &  & Визначено частоту або ситуацію, що вимагає проведення аудиту для кожної ідентифікованої події &  & Параметр: au\_\allowbreak{}2\_\allowbreak{}odp\_\allowbreak{}03\newline Тип: list\newline Значення: \textbf{login, logout, failed\_\allowbreak{}attempt} \\ \cline{3-5}
 &  & Частота перегляду та оновлення типів подій, обраних для журналювання &  & Параметр: au\_\allowbreak{}2\_\allowbreak{}odp\_\allowbreak{}04\newline Тип: string\newline Значення: \textbf{щорічно} \\ \hline
20 & ЗМІСТ ЗАПИСІВ АУДИТУ (AU-3) & Записи аудиту містять інформацію, яка встановлює який тип події стався & AU-3 & Параметр: au\_\allowbreak{}3\_\allowbreak{}a\newline Тип: list\newline Значення: \textbf{login, logout, failed\_\allowbreak{}attempt} \\ \cline{3-5}
 &  & Записи аудиту містять інформацію, яка встановлює джерело події &  & Параметр: au\_\allowbreak{}3\_\allowbreak{}d\newline Тип: list\newline Значення: \textbf{login, logout, failed\_\allowbreak{}attempt} \\ \cline{3-5}
 &  & Записи аудиту містять інформацію, яка встановлює наслідки події &  & Параметр: au\_\allowbreak{}3\_\allowbreak{}e\newline Тип: list\newline Значення: \textbf{login, logout, failed\_\allowbreak{}attempt} \\ \cline{3-5}
 &  & Записи аудиту містять інформацію, яка встановлює результат події та ідентифікатор будь-яких осіб або суб'єктів, пов'язаних з подією &  & Параметр: au\_\allowbreak{}3\_\allowbreak{}f\newline Тип: list\newline Значення: \textbf{login, logout, failed\_\allowbreak{}attempt} \\ \cline{3-5}
 &  & & AU-3(1) & \\ \hline
21 & ЗБЕРЕЖЕННЯ ЗАПИСІВ АУДИТУ (AU-11) & & AU-11 & \\ \cline{3-5}
 &  & визначено персонал або ролі, яким дозволено обирати типи подій, що мають реєструватися певними компонентами системи; & AU-12 & Параметр: au\_\allowbreak{}12\_\allowbreak{}odp\_\allowbreak{}2\newline Тип: list\newline Значення: \textbf{admin} \\ \cline{3-5}
 &  & <AU-12\_\allowbreak{}ODP[02] персонал або ролі> може/можуть вибирати типи подій, які будуть реєструватися певними компонентами системи; &  & Параметр: au\_\allowbreak{}12\_\allowbreak{}b\newline Тип: list\newline Значення: \textbf{admin} \\ \hline
22 & РЕАГУВАННЯ НА ВІДМОВИ ОБРОБКИ ДАНИХ АУДИТУ (AU-5) & Персонал або ролі отримують сповіщення у разі збою процесу обробки даних аудиту періоду часу & AU-5 & Параметр: au\_\allowbreak{}5\_\allowbreak{}a\newline Тип: list\newline Значення: \textbf{admin, security\_\allowbreak{}officer} \\ \cline{3-5}
 &  & Додаткові дії виконуються у разі збою процесу обробки даних аудиту &  & Параметр: au\_\allowbreak{}5\_\allowbreak{}b\newline Тип: list\newline Значення: \textbf{login, logout, failed\_\allowbreak{}attempt} \\ \cline{3-5}
 &  & Визначено персонал або ролі, які отримують сповіщення про збої в процесі обробки даних аудиту  &  & Параметр: au\_\allowbreak{}5\_\allowbreak{}odp\_\allowbreak{}01\newline Тип: list\newline Значення: \textbf{admin, security\_\allowbreak{}officer} \\ \cline{3-5}
 &  & Визначено період часу, протягом якого персонал або ролі отримують сповіщення про збої в процесі обробки даних аудиту &  & Параметр: au\_\allowbreak{}5\_\allowbreak{}odp\_\allowbreak{}02\newline Тип: list\newline Значення: \textbf{admin, security\_\allowbreak{}officer} \\ \cline{3-5}
 &  & Визначено додаткові дії, яких слід вжити у випадку збою в процесі обробки даних аудиту  &  & Параметр: au\_\allowbreak{}5\_\allowbreak{}odp\_\allowbreak{}03\newline Тип: list\newline Значення: \textbf{login, logout, failed\_\allowbreak{}attempt} \\ \hline
23 & ОГЛЯД, АНАЛІЗ І ЗВІТНІСТЬ АУДИТУ (AU-6) & Записи аудиту системи переглядаються та аналізуються частота для виявлення ознак неналежної або незвичної діяльності та потенційного впливу неналежної або незвичної діяльності & AU-6 & Параметр: au\_\allowbreak{}6\_\allowbreak{}a\newline Тип: string\newline Значення: \textbf{щотижня} \\ \cline{3-5}
 &  & Звіт аудиту відправляється персоналу або ролям &  & Параметр: au\_\allowbreak{}6\_\allowbreak{}b\newline Тип: list\newline Значення: \textbf{admin, security\_\allowbreak{}officer} \\ \cline{3-5}
 &  & Визначено частоту, з якою переглядаються та аналізуються записи аудиту системи &  & Параметр: au\_\allowbreak{}6\_\allowbreak{}odp\_\allowbreak{}01\newline Тип: integer\newline Значення: \textbf{30} \\ \cline{3-5}
 &  & Визначено персонал або ролі які отримують результати оглядів та аналізів системних записів &  & Параметр: au\_\allowbreak{}6\_\allowbreak{}odp\_\allowbreak{}03\newline Тип: list\newline Значення: \textbf{admin, security\_\allowbreak{}officer} \\ \cline{3-5}
 &  & & AU-6(3) & \\ \hline
24 & СКОРОЧЕННЯ ЗАПИСІВ АУДИТУ ТА ФОРМУВАННЯ ЗВІТУ (AU-7) & & AU-7 & \\ \hline
25 & ПОЗНАЧКА ЧАСУ (AU-8) & & AU-8 & \\ \hline
26 & ЗАХИСТ ІНФОРМАЦІЇ АУДИТУ (AU-9) & визначено персонал або ролі, які мають бути сповіщені при виявленні несанкціонованого доступу, зміни або видалення інформації аудиту; & AU-9 & Параметр: au\_\allowbreak{}9\_\allowbreak{}odp\newline Тип: list\newline Значення: \textbf{admin} \\ \cline{3-5}
 &  & <AU-09\_\allowbreak{}ODP персонал або ролі>  отримують сповіщення при виявленні несанкціонованого доступу, зміни або видалення інформації аудиту. &  & Параметр: au\_\allowbreak{}9\_\allowbreak{}b\newline Тип: list\newline Значення: \textbf{admin} \\ \cline{3-5}
 &  & & AU-9(4) & \\ \hline
\multicolumn{5}{|>{\raggedright\arraybackslash}p{16cm}|}{\textbf{УПРАВЛІННЯ КОНФІГУРАЦІЄЮ (CM)}} \\* \noalign{\hrule height 0.4pt}
27 & БАЗОВА КОНФІГУРАЦІЯ (CM-2) & & CM-2 & \\ \hline
28 & НАЛАШТУВАННЯ КОНФІГУРАЦІЇ (CM-6) & & CM-6 & \\ \hline
29 & УПРАВЛІННЯ ЗМІНАМИ КОНФІГУРАЦІЇ (CM-3) & & CM-3 & \\ \hline
30 & АНАЛІЗ ВПЛИВУ НА БЕЗПЕКУ ТА ПРИВАТНІСТЬ (CM-4) & & CM-4 & \\ \cline{3-5}
 &  & & CM-4(2) & \\ \hline
31 & ОБМЕЖЕННЯ ДОСТУПУ ДО ЗМІНИ (CM-5) & & CM-5 & \\ \hline
32 & МІНІМАЛЬНО НЕОБХІДНА ФУНКЦІОНАЛЬНІСТЬ (CM-7) & & CM-7 & \\ \cline{3-5}
 &  & & CM-7(1) & \\ \hline
33 & АВТОРИЗОВАНЕ ПРОГРАМНЕ ЗАБЕЗПЕЧЕННЯ -- БІЛИЙ СПИСОК (CM-7(5)) (CM-7) & & CM-7(5) & \\ \hline
34 & ІНВЕНТАРИЗАЦІЯ КОМПОНЕНТІВ СИСТЕМИ (CM-8) & & CM-8 & \\ \cline{3-5}
 &  & & CM-8(1) & \\ \hline
35 & РОЗТАШУВАННЯ ІНФОРМАЦІЇ (CM-12) & & CM-12 & \\ \hline
36 & КОНФІГУРАЦІЯ СИСТЕМ ТА КОМПОНЕНТІВ ДЛЯ СФЕР З ВИСОКИМ РИЗИКОМ (CM-2(7)) (CM-2) & & CM-2(7) & \\ \hline
\multicolumn{5}{|>{\raggedright\arraybackslash}p{16cm}|}{\textbf{ІДЕНТИФІКАЦІЯ ТА АВТЕНТИФІКАЦІЯ (IA)}} \\* \noalign{\hrule height 0.4pt}
37 & ІДЕНТИФІКАЦІЯ ТА АВТЕНТИФІКАЦІЯ КОРИСТУВАЧІВ (IA-2) & & IA-2 & \\ \cline{3-5}
 &  & & IA-11 & \\ \hline
38 & ІДЕНТИФІКАЦІЯ ТА АВТЕНТИФІКАЦІЯ ПРИСТРОЇВ (IA-3) & & IA-3 & \\ \hline
39 & БАГАТОФАКТОРНА АВТЕНТИФІКАЦІЯ ПРИВІЛЕЙОВАНИХ ОБЛІКОВИХ ЗАПИСІВ (IA-2(1)) (IA-2) & & IA-2(1) & \\ \cline{3-5}
 &  & & IA-2(2) & \\ \hline
40 & ДОСТУП ДО ОБЛІКОВИХ ЗАПИСІВ --- СТІЙКІСТЬ ДО ВІДТВОРЕННЯ (IA-2(8)) (IA-2) & & IA-2(8) & \\ \hline
41 & УПРАВЛІННЯ ІДЕНТИФІКАЦІЄЮ (IA-4) & & IA-4 & \\ \cline{3-5}
 &  & & IA-4(4) & \\ \hline
42 & АВТЕНТИФІКАЦІЯ НА ОСНОВІ ПАРОЛЯ (IA-5(1)) (IA-5) & & IA-5(1) & \\ \hline
43 & ПРИХОВУВАННЯ ЗВОРОТНОГО ЗВ'ЯЗКУ АВТЕНТИФІКАТОРА (IA-6) & & IA-6 & \\ \hline
44 & АВТЕНТИФІКАТОР УПРАВЛІННЯ (IA-5) & & IA-5 & \\ \hline
\multicolumn{5}{|>{\raggedright\arraybackslash}p{16cm}|}{\textbf{РЕАГУВАННЯ НА ІНЦИДЕНТИ (IR)}} \\* \noalign{\hrule height 0.4pt}
45 & ОБРОБКА ІНЦИДЕНТУ (IR-4) & & IR-4 & \\ \hline
46 & МОНІТОРИНГ ІНЦИДЕНТУ (IR-5) & & IR-5 & \\ \cline{3-5}
 &  & & IR-6 & \\ \cline{3-5}
 &  & & IR-7 & \\ \hline
47 & ПЕРЕВІРКА РЕАГУВАНЬ НА ІНЦИДЕНТИ (IR-3) & & IR-3 & \\ \hline
48 & НАВЧАННЯ З РЕАГУВАННЯ НА ІНЦИДЕНТИ (IR-2) & & IR-2 & \\ \hline
49 & ПЛАН РЕАГУВАННЯ НА ІНЦИДЕНТИ (IR-8) & & IR-8 & \\ \hline
\multicolumn{5}{|>{\raggedright\arraybackslash}p{16cm}|}{\textbf{ТЕХНІЧНЕ ОБСЛУГОВУВАННЯ (MA)}} \\* \noalign{\hrule height 0.4pt}
50 & ІНСТРУМЕНТИ ДЛЯ ОБСЛУГОВУВАННЯ (MA-3) & & MA-3 & \\ \cline{3-5}
 &  & & MA-3(1) & \\ \cline{3-5}
 &  & & MA-3(2) & \\ \cline{3-5}
 &  & & MA-3(3) & \\ \hline
51 & ВІДДАЛЕНЕ ОБСЛУГОВУВАННЯ (MA-4) & & MA-4 & \\ \hline
52 & ТЕХНІЧНИЙ ПЕРСОНАЛ (MA-5) & & MA-5 & \\ \hline
\multicolumn{5}{|>{\raggedright\arraybackslash}p{16cm}|}{\textbf{ЗАХИСТ НОСІЇВ ІНФОРМАЦІЇ (MP)}} \\* \noalign{\hrule height 0.4pt}
53 & ЗБЕРІГАННЯ НОСІЇВ ІНФОРМАЦІЇ (MP-4) & & MP-4 & \\ \hline
54 & ДОСТУП ДО НОСІЇВ ІНФОРМАЦІЇ (MP-2) & & MP-2 & \\ \hline
55 & ЗНИЩЕННЯ ІНФОРМАЦІЇ НА НОСІЯХ ІНФОРМАЦІЇ (MP-6) & & MP-6 & \\ \hline
56 & МАРКУВАННЯ НОСІЇВ ІНФОРМАЦІЇ (MP-3) & & MP-3 & \\ \hline
57 & ТРАНСПОРТУВАННЯ НОСІЇВ ІНФОРМАЦІЇ (MP-5) & & MP-5 & \\ \cline{3-5}
 &  & & SC-28 & \\ \hline
58 & ВИКОРИСТАННЯ НОСІЇВ ІНФОРМАЦІЇ (MP-7) & & MP-7 & \\ \hline
\multicolumn{5}{|>{\raggedright\arraybackslash}p{16cm}|}{\textbf{ПЛАНУВАННЯ БЕЗПЕРЕРВНОЇ РОБОТИ (CP)}} \\* \noalign{\hrule height 0.4pt}
59 & РЕЗЕРВНЕ КОПІЮВАННЯ (CP-9) & & CP-9 & \\ \cline{3-5}
 &  & & CP-9(8) & \\ \hline
\multicolumn{5}{|>{\raggedright\arraybackslash}p{16cm}|}{\textbf{БЕЗПЕКА ПЕРСОНАЛУ (PS)}} \\* \noalign{\hrule height 0.4pt}
60 & ПЕРЕВІРКА ПЕРСОНАЛУ (PS-3) & & PS-3 & \\ \hline
61 & ЗВІЛЬНЕННЯ ПЕРСОНАЛУ (PS-4) & & PS-4 & \\ \cline{3-5}
 &  & & PS-5 & \\ \hline
\multicolumn{5}{|>{\raggedright\arraybackslash}p{16cm}|}{\textbf{ФІЗИЧНИЙ ЗАХИСТ ТА ЗАХИСТ НАВКОЛИШНЬОГО СЕРЕДОВИЩА (PE)}} \\* \noalign{\hrule height 0.4pt}
62 & РЕ-2 ДОСТУПУ (PE-2) & & PE-2 & \\ \hline
63 & МОНІТОРИНГ ФІЗИЧНОГО ДОСТУПУ (PE-6) & & PE-6 & \\ \hline
64 & АЛЬТЕРНАТИВНЕ РОБОЧЕ МІСЦЕ (PE-17) & & PE-17 & \\ \hline
65 & ФІЗИЧНИЙ ДОСТУП ДО СИСТЕМИ (PE-3) & & PE-3 & \\ \cline{3-5}
 &  & & PE-5 & \\ \hline
66 & РЕ-4 ЛІНІЙ ЕЛЕКТРОЖИВЛЕННЯ (PE-4) & & PE-4 & \\ \hline
\multicolumn{5}{|>{\raggedright\arraybackslash}p{16cm}|}{\textbf{ОЦІНКА РИЗИКІВ (RA)}} \\* \noalign{\hrule height 0.4pt}
67 & ОЦІНЮВАННЯ РИЗИКУ (RA-3) & & RA-3 & \\ \cline{3-5}
 &  & & RA-3(1) & \\ \cline{3-5}
 &  & & SR-6 & \\ \hline
68 & СКАНУВАННЯ ВРАЗЛИВОСТЕЙ (RA-5) & & RA-5 & \\ \cline{3-5}
 &  & & RA-5(2) & \\ \hline
69 & РЕАГУВАННЯ НА РИЗИК (RA-7) & & RA-7 & \\ \hline
\multicolumn{5}{|>{\raggedright\arraybackslash}p{16cm}|}{\textbf{ОЦІНЮВАННЯ, АВТОРИЗАЦІЯ ТА МОНІТОРИНГ (CA)}} \\* \noalign{\hrule height 0.4pt}
70 & ОЦІНЮВАННЯ (CA-2) & визначено частоту, з якою слід оцінювати засоби контролю в системі та середовищі її функціонування; & CA-2 & Параметр: ca\_\allowbreak{}2\_\allowbreak{}odp\_\allowbreak{}1\newline Тип: list\newline Значення: \textbf{admin} \\ \cline{3-5}
 &  & визначені особи або ролі, яким мають бути надані результати оцінювання з безпеки; &  & Параметр: ca\_\allowbreak{}2\_\allowbreak{}odp\_\allowbreak{}2\newline Тип: list\newline Значення: \textbf{admin} \\ \cline{3-5}
 &  & розроблено план контрольної оцінки, який описує обсяг оцін ки, включаючи заходи захисту та посилені заходи, що підлягають оцінюванню. &  & Параметр: ca\_\allowbreak{}2\_\allowbreak{}b\_\allowbreak{}1\newline Тип: list\newline Значення: \textbf{admin} \\ \cline{3-5}
 &  & розроблено план контрольної оцінки, який  описує обсяг оцінки, включаючи процедури оцінювання, які використовуватимуться для визначення ефективності заходів. &  & Параметр: ca\_\allowbreak{}2\_\allowbreak{}b\_\allowbreak{}2\newline Тип: list\newline Значення: \textbf{admin} \\ \cline{3-5}
 &  & розроблено план контрольної оцінки, який описує обсяг оцінки, включаючи групу оцінювання. &  & Параметр: ca\_\allowbreak{}2\_\allowbreak{}b\_\allowbreak{}3\_\allowbreak{}2\newline Тип: list\newline Значення: \textbf{admin} \\ \cline{3-5}
 &  & розроблено план контрольної оцінки, який описує обсяг оцінки, включаючи ролі й обов'язки з оцінювання. &  & Параметр: ca\_\allowbreak{}2\_\allowbreak{}b\_\allowbreak{}3\_\allowbreak{}3\newline Тип: list\newline Значення: \textbf{admin} \\ \cline{3-5}
 &  & результати оцінювання з безпеки надаються <CA-02\_\allowbreak{}ODP[02] особам або ролям>. &  & Параметр: ca\_\allowbreak{}2\_\allowbreak{}f\newline Тип: list\newline Значення: \textbf{admin} \\ \hline
71 & ПЛАН УСУНЕННЯ НЕДОЛІКІВ ТА КОНТРОЛЬНІ ПОКАЗНИКИ (CA-5) & визначено частоту оновлення чинного плану усунення недоліків та контрольних показників на основі результатів оцінювання заходів захисту, незалежних аудитів та постійного моніторингу; & CA-5 & Параметр: ca\_\allowbreak{}5\_\allowbreak{}odp\newline Тип: list\newline Значення: \textbf{admin} \\ \cline{3-5}
 &  & розроблено план усунення недоліків та контрольні показники для системи, та задокументувано заплановані дії організації з коригування, спрямовані на усунення недоліків та зауважень, виявлених під час оцінки заходів захисту, а також на зменшення або усунення відомих вразливостей в системі; &  & Параметр: ca\_\allowbreak{}5\_\allowbreak{}a\newline Тип: list\newline Значення: \textbf{admin} \\ \hline
72 & БЕЗПЕРЕРВНИЙ МОНІТОРИНГ (CA-7) & визначено частоту, з якою слід моніторити ефективність заходів захисту; & CA-7 & Параметр: ca\_\allowbreak{}7\_\allowbreak{}odp\_\allowbreak{}2\newline Тип: integer\newline Значення: \textbf{30} \\ \cline{3-5}
 &  & визначено частоту, з якою слід оцінювати ефективність заходів захисту; &  & Параметр: ca\_\allowbreak{}7\_\allowbreak{}odp\_\allowbreak{}3\newline Тип: integer\newline Значення: \textbf{30} \\ \cline{3-5}
 &  & визначено персонал або ролі, яким повідомляється про стан безпеки системи; &  & Параметр: ca\_\allowbreak{}7\_\allowbreak{}odp\_\allowbreak{}4\newline Тип: list\newline Значення: \textbf{admin} \\ \cline{3-5}
 &  & визначено частоту, з якою повідомляється про стан безпеки системи; &  & Параметр: ca\_\allowbreak{}7\_\allowbreak{}odp\_\allowbreak{}5\newline Тип: integer\newline Значення: \textbf{30} \\ \cline{3-5}
 &  & визначено персонал або ролі, яким повідомляється про стан конфіденційності системи; &  & Параметр: ca\_\allowbreak{}7\_\allowbreak{}odp\_\allowbreak{}6\newline Тип: list\newline Значення: \textbf{admin} \\ \cline{3-5}
 &  & визначено частоту, з якою повідомляється про стан конфіденційності системи; &  & Параметр: ca\_\allowbreak{}7\_\allowbreak{}odp\_\allowbreak{}7\newline Тип: integer\newline Значення: \textbf{30} \\ \cline{3-5}
 &  & безперервний моніторинг на рівні системи включає поточні контрольні оцінки відповідно до стратегії безперервного моніторингу; &  & Параметр: ca\_\allowbreak{}7\_\allowbreak{}c\newline Тип: list\newline Значення: \textbf{admin} \\ \hline
73 & ВЗАЄМОДІЯ СИСТЕМ (CA-3) & вибрано одне або декілька з наступних ЗНАЧЕНЬ ПАРАМЕТРІВ: \{): угоди безпеки взаємозв'язку; & CA-3 & Параметр: ca\_\allowbreak{}3\_\allowbreak{}odp\_\allowbreak{}1\newline Тип: list\newline Значення: \textbf{admin} \\ 
 &  & договори безпеки обміну інформацією; &  &  \\ 
 &  & меморандуми про взаєморозуміння; &  &  \\ 
 &  & угоди про рівень обслуговування; &  &  \\ 
 &  & угоди користувача; &  &  \\ 
 &  & угоди &  &  \\ \cline{3-5}
 &  & визначено частоту, з якою необхідно переглядати та оновлювати угоди; &  & Параметр: ca\_\allowbreak{}3\_\allowbreak{}odp\_\allowbreak{}3\newline Тип: integer\newline Значення: \textbf{30} \\ \hline
\multicolumn{5}{|>{\raggedright\arraybackslash}p{16cm}|}{\textbf{ЗАХИСТ СИСТЕМ ТА КОМУНІКАЦІЙ (SC)}} \\* \noalign{\hrule height 0.4pt}
74 & ДОСТУПНІСТЬ РЕСУРСІВ (SC-7) & & SC-7 & \\ \hline
75 & ІНФОРМАЦІЯ В ЗАГАЛЬНИХ СИСТЕМНИХ РЕСУРСАХ (SC-4) & & SC-4 & \\ \hline
76 & ВІДМОВА ЗА ЗАМОВЧУВАННЯМ - ДОЗВІЛ ЗА ВИНЯТКОМ (SC-7(5)) (SC-7) & & SC-7(5) & \\ \hline
77 & КОНФІДЕНЦІЙНІСТЬ ТА ЦІЛІСНІСТЬ ПЕРЕДАЧІ (SC-8) & & SC-8 & \\ \cline{3-5}
 &  & & SC-8(1) & \\ \cline{3-5}
 &  & & SC-28 & \\ \cline{3-5}
 &  & Реалізовані криптографічні механізми для запобігання несанкціонованому розкриттю інформації, що знаходиться в стані спокою на системних компонентах або носіях & SC-28(1) & Параметр: sc\_\allowbreak{}28\_\allowbreak{}1\_\allowbreak{}01\newline Тип: string\newline Значення: \textbf{AES-256-GCM} \\ \cline{3-5}
 &  & Реалізовані криптографічні механізми для запобігання несанкціонованій модифікації інформації, що знаходиться в стані спокою на системних компонентах або носіях &  & Параметр: sc\_\allowbreak{}28\_\allowbreak{}1\_\allowbreak{}02\newline Тип: string\newline Значення: \textbf{AES-256-GCM} \\ \hline
78 & ВІДКЛЮЧЕННЯ МЕРЕЖІ (SC-10) & Визначено період бездіяльності, після якого система розриває мережеве з'єднання, пов'язане з сеансом зв'язку; & SC-10 & Параметр: sc\_\allowbreak{}10\_\allowbreak{}odp\newline Тип: integer\newline Значення: \textbf{30} \\ 
 &  & SC08(05)\_\allowbreak{}ODP[02] мережеве з'єднання, пов'язане з сеансом зв'язку, розірвано в кінці сеансу або після період часу бездіяльності &  &  \\ \hline
79 & ВСТАНОВЛЕННЯ КЛЮЧАМИ (SC-12) & Встановлюються криптографічні ключі, коли в системі використовується криптографія відповідно до < SC-12\_\allowbreak{}ODP вимог > & SC-12 & Параметр: sc\_\allowbreak{}12\_\allowbreak{}01\newline Тип: string\newline Значення: \textbf{AES-256-GCM} \\ \cline{3-5}
 &  & Здійснюється управління криптографічними ключами, коли в системі використовується криптографія, відповідно до вимог  &  & Параметр: sc\_\allowbreak{}12\_\allowbreak{}02\newline Тип: string\newline Значення: \textbf{AES-256-GCM} \\ \hline
80 & КРИПТОГРАФІЧНИЙ ЗАХИСТ (SC-13) & КРИПТОГРАФІЧНИЙ ЗАХИСТ МЕТА ОЦІНКИ: Визначити, чи: & SC-13 & Параметр: sc\_\allowbreak{}13\_\allowbreak{}01\newline Тип: string\newline Значення: \textbf{AES-256-GCM} \\ \cline{3-5}
 &  & Визначено використання криптографічних засобів &  & Параметр: sc\_\allowbreak{}13\_\allowbreak{}odp\_\allowbreak{}01\newline Тип: string\newline Значення: \textbf{AES-256-GCM} \\ \cline{3-5}
 &  & Визначено типи криптографії для кожного вказаного криптографічного використання; &  & Параметр: sc\_\allowbreak{}13\_\allowbreak{}odp\_\allowbreak{}02\newline Тип: string\newline Значення: \textbf{AES-256-GCM} \\ 
 &  & SC-13a. ідентифіковано використання>; &  &  \\ 
 &  & SC-13b. реалізовано типи криптографії для кожного вказаного криптографічного використання (визначеного в SC-13\_\allowbreak{}ODP[01]). криптографічне <SC-13\_\allowbreak{}ODP[01] &  &  \\ \hline
81 & СПІЛЬНІ ОБЧИСЛЮВАЛЬНІ ПРИСТРОЇ ТА ЗАСТОСУНКИ (SC-15) & & SC-15 & \\ \hline
82 & МОБІЛЬНИЙ КОД (SC-18) & & SC-18 & \\ \hline
83 & АВТЕНТИФІКАЦІЯ СЕСІЇ (SC-23) & & SC-23 & \\ \hline
\multicolumn{5}{|>{\raggedright\arraybackslash}p{16cm}|}{\textbf{ЦІЛІСНІСТЬ СИСТЕМИ ТА ІНФОРМАЦІЇ (SI)}} \\* \noalign{\hrule height 0.4pt}
84 & ВИПРАВЛЕННЯ ДЕФЕКТІВ (SI-2) & & SI-2 & \\ \hline
85 & ЗАХИСТ ВІД ШКІДЛИВОГО КОДУ (SI-3) & Визначено частоту, з якою механізми шкідливого коду виконують сканування & SI-3 & Параметр: si\_\allowbreak{}3\_\allowbreak{}odp\_\allowbreak{}02\newline Тип: integer\newline Значення: \textbf{30} \\ \cline{3-5}
 &  & Визначено дії, яких слід вжити у відповідь на виявлення шкідливого коду (якщо вибрано) &  & Параметр: si\_\allowbreak{}3\_\allowbreak{}odp\_\allowbreak{}05\newline Тип: list\newline Значення: \textbf{login, logout, failed\_\allowbreak{}attempt} \\ \hline
86 & ПОПЕРЕДЖЕННЯ, РЕКОМЕНДАЦІЇ ТА ДИРЕКТИВИ З БЕЗПЕКИ (SI-5) & Вибрано одне або декілька з наступних ЗНАЧЕНЬ ПАРАМЕТРІВ: \{персонал або ролі; & SI-5 & Параметр: si\_\allowbreak{}5\_\allowbreak{}odp\_\allowbreak{}02\newline Тип: list\newline Значення: \textbf{admin, security\_\allowbreak{}officer} \\ 
 &  & елементи; &  &  \\ 
 &  & зовнішні організації\} &  &  \\ \cline{3-5}
 &  & Визначено персонал або ролі, до яких мають бути доведені попередження, поради та директиви з безпеки (якщо визначено) &  & Параметр: si\_\allowbreak{}5\_\allowbreak{}odp\_\allowbreak{}03\newline Тип: list\newline Значення: \textbf{admin, security\_\allowbreak{}officer} \\ \hline
87 & МОНІТОРИНГ СИСТЕМИ (SI-4) & Визначені методи та способи, що використовуються для виявлення несанкціонованого використання системи & SI-4 & Параметр: si\_\allowbreak{}4\_\allowbreak{}odp\_\allowbreak{}02\newline Тип: list\newline Значення: \textbf{admin, security\_\allowbreak{}officer} \\ \cline{3-5}
 &  & Визначена інформація про моніторинг системи, яка повинна надаватися персоналу або ролям &  & Параметр: si\_\allowbreak{}4\_\allowbreak{}odp\_\allowbreak{}03\newline Тип: list\newline Значення: \textbf{admin, security\_\allowbreak{}officer} \\ \cline{3-5}
 &  & Визначено персонал або ролі, яким має надаватися інформація про моніторинг системи &  & Параметр: si\_\allowbreak{}4\_\allowbreak{}odp\_\allowbreak{}04\newline Тип: list\newline Значення: \textbf{admin, security\_\allowbreak{}officer} \\ \cline{3-5}
 &  & Вибрано одне або декілька з наступних ЗНАЧЕНЬ ПАРАМЕТРІВ: \{за потребою; &  & Параметр: si\_\allowbreak{}4\_\allowbreak{}odp\_\allowbreak{}05\newline Тип: string\newline Значення: \textbf{щорічно} \\ 
 &  & частота\} &  &  \\ \cline{3-5}
 &  & Визначені критерії незвичної або несанкціонованої діяльності або умови для вхідного трафіку зв'язку & SI-4(4) & Параметр: si\_\allowbreak{}4\_\allowbreak{}4\_\allowbreak{}a\_\allowbreak{}01\newline Тип: list\newline Значення: \textbf{} \\ \cline{3-5}
 &  & Визначені критерії незвичайної або несанкціонованої діяльності або умови для вихідного трафіку зв'язку &  & Параметр: si\_\allowbreak{}4\_\allowbreak{}4\_\allowbreak{}a\_\allowbreak{}02\newline Тип: list\newline Значення: \textbf{} \\ \cline{3-5}
 &  & Здійснюється моніторинг вхідного комунікаційного трафіку частота на предмет незвичних або несанкціонованих дій або умов &  & Параметр: si\_\allowbreak{}4\_\allowbreak{}4\_\allowbreak{}b\_\allowbreak{}01\newline Тип: string\newline Значення: \textbf{щорічно} \\ \cline{3-5}
 &  & Контролюється вихідний трафік зв'язку частота на предмет незвичних або несанкціонованих дій або умов. вихідного виявлення &  & Параметр: si\_\allowbreak{}4\_\allowbreak{}4\_\allowbreak{}b\_\allowbreak{}02\newline Тип: string\newline Значення: \textbf{щорічно} \\ \hline
88 & УПРАВЛІННЯ ТА ЗБЕРЕЖЕННЯ ІНФОРМАЦІЇ (SI-12) & Здійснюється управління інформацією в системі відповідно до чинних законів, наказів, директив, положень, політик, стандартів, інструкцій та експлуатаційних вимог & SI-12 & Параметр: si\_\allowbreak{}12\_\allowbreak{}01\newline Тип: list\newline Значення: \textbf{default\_\allowbreak{}deny\_\allowbreak{}rule, abac\_\allowbreak{}rule\_\allowbreak{}1} \\ \cline{3-5}
 &  & Зберігається інформація в системі відповідно до чинних законів, указів Президента, директив, положень, політик, стандартів, інструкцій та експлуатаційних вимог &  & Параметр: si\_\allowbreak{}12\_\allowbreak{}02\newline Тип: list\newline Значення: \textbf{default\_\allowbreak{}deny\_\allowbreak{}rule, abac\_\allowbreak{}rule\_\allowbreak{}1} \\ \cline{3-5}
 &  & Управління інформацією, що виводиться з системи, здійснюється відповідно до чинних законів, указів Президента, директив, положень, політик, стандартів, інструкцій та операційних вимог &  & Параметр: si\_\allowbreak{}12\_\allowbreak{}03\newline Тип: list\newline Значення: \textbf{default\_\allowbreak{}deny\_\allowbreak{}rule, abac\_\allowbreak{}rule\_\allowbreak{}1} \\ \cline{3-5}
 &  & Зберігається інформація, що виводиться з системи, відповідно до чинних законів, наказів, директив, положень, політик, стандартів, інструкцій та експлуатаційних вимог &  & Параметр: si\_\allowbreak{}12\_\allowbreak{}04\newline Тип: list\newline Значення: \textbf{default\_\allowbreak{}deny\_\allowbreak{}rule, abac\_\allowbreak{}rule\_\allowbreak{}1} \\ \hline
\multicolumn{5}{|>{\raggedright\arraybackslash}p{16cm}|}{\textbf{Політики та процедури з безпеки}} \\* \noalign{\hrule height 0.4pt}
89 & Політики та процедури з безпеки & Визначено персонал або ролі, на які поширюється політика контролю доступу & AC-1 & Параметр: ac\_\allowbreak{}1\_\allowbreak{}odp\_\allowbreak{}01\newline Тип: list\newline Значення: \textbf{admin, security\_\allowbreak{}officer} \\ \cline{3-5}
 &  & Визначено персонал або ролі, на які поширюються процедури контролю доступу &  & Параметр: ac\_\allowbreak{}1\_\allowbreak{}odp\_\allowbreak{}02\newline Тип: list\newline Значення: \textbf{admin, security\_\allowbreak{}officer} \\ \cline{3-5}
 &  & Визначено посадову особу, яка керуватиме політикою та процедурами контролю доступу &  & Параметр: ac\_\allowbreak{}1\_\allowbreak{}odp\_\allowbreak{}04\newline Тип: list\newline Значення: \textbf{admin, security\_\allowbreak{}officer} \\ \cline{3-5}
 &  & Визначено частоту, з якою переглядається та оновлюється поточна політика контролю доступу &  & Параметр: ac\_\allowbreak{}1\_\allowbreak{}odp\_\allowbreak{}05\newline Тип: list\newline Значення: \textbf{default\_\allowbreak{}deny\_\allowbreak{}rule, abac\_\allowbreak{}rule\_\allowbreak{}1} \\ \cline{3-5}
 &  & Визначено події, які вимагають перегляду та оновлення поточної політики контролю доступу &  & Параметр: ac\_\allowbreak{}1\_\allowbreak{}odp\_\allowbreak{}06\newline Тип: list\newline Значення: \textbf{default\_\allowbreak{}deny\_\allowbreak{}rule, abac\_\allowbreak{}rule\_\allowbreak{}1} \\ \cline{3-5}
 &  & Визначено частоту, з якою переглядаються та оновлюються поточні процедури контролю доступу &  & Параметр: ac\_\allowbreak{}1\_\allowbreak{}odp\_\allowbreak{}07\newline Тип: integer\newline Значення: \textbf{30} \\ \cline{3-5}
 &  & Визначено події, які потребують перегляду та оновлення процедур &  & Параметр: ac\_\allowbreak{}1\_\allowbreak{}odp\_\allowbreak{}08\newline Тип: integer\newline Значення: \textbf{30} \\ \cline{3-5}
 &  & Розроблено та задокументовано політику контролю доступу &  & Параметр: ac\_\allowbreak{}1\_\allowbreak{}a\_\allowbreak{}01\newline Тип: integer\newline Значення: \textbf{30} \\ \cline{3-5}
 &  & Політика контролю доступу поширюється на <AC-01\_\allowbreak{}ODP[01] персонал або ролі>; &  & Параметр: ac\_\allowbreak{}1\_\allowbreak{}a\_\allowbreak{}02\newline Тип: integer\newline Значення: \textbf{30} \\ \cline{3-5}
 &  & Розроблені та задокументовані процедури контролю доступу для полегшення впровадження політики контролю доступу та пов'язаних з нею заходів захисту &  & Параметр: ac\_\allowbreak{}1\_\allowbreak{}a\_\allowbreak{}03\newline Тип: integer\newline Значення: \textbf{30} \\ \cline{3-5}
 &  & Процедури контролю доступу поширюються на <AC-01\_\allowbreak{}ODP[02] персонал або ролі> &  & Параметр: ac\_\allowbreak{}1\_\allowbreak{}a\_\allowbreak{}04\newline Тип: integer\newline Значення: \textbf{30} \\ \cline{3-5}
 &  & Політика контролю доступу <AC-01\_\allowbreak{}ODP[03] ВИБІРКОВЕ ЗНАЧЕННЯ ПАРАМЕТРА(ів)> містить мету &  & Параметр: ac\_\allowbreak{}1\_\allowbreak{}a\_\allowbreak{}1\_\allowbreak{}a\_\allowbreak{}01\newline Тип: integer\newline Значення: \textbf{30} \\ \cline{3-5}
 &  & Політика контролю доступу <AC-01\_\allowbreak{}ODP[03] ВИБІРКОВЕ ЗНАЧЕННЯ ПАРАМЕТРА(ів)> містить сферу застосування &  & Параметр: ac\_\allowbreak{}1\_\allowbreak{}a\_\allowbreak{}1\_\allowbreak{}a\_\allowbreak{}02\newline Тип: integer\newline Значення: \textbf{30} \\ \cline{3-5}
 &  & Політика контролю доступу <AC-01\_\allowbreak{}ODP[03] ВИБІРКОВЕ ЗНАЧЕННЯ ПАРАМЕТРА(ів)> містить сферу застосування &  & Параметр: ac\_\allowbreak{}1\_\allowbreak{}a\_\allowbreak{}1\_\allowbreak{}a\_\allowbreak{}04\newline Тип: integer\newline Значення: \textbf{30} \\ \cline{3-5}
 &  & Політика контролю доступу <AC-01\_\allowbreak{}ODP[03] ВИБІРКОВЕ ЗНАЧЕННЯ ПАРАМЕТРА(ів)> містить сферу застосування &  & Параметр: ac\_\allowbreak{}1\_\allowbreak{}a\_\allowbreak{}1\_\allowbreak{}a\_\allowbreak{}05\newline Тип: integer\newline Значення: \textbf{30} \\ \cline{3-5}
 &  & Політика контролю доступу <AC-01\_\allowbreak{}ODP[03] ВИБІРКОВЕ ЗНАЧЕННЯ ПАРАМЕТРА(ів)> містить сферу застосування &  & Параметр: ac\_\allowbreak{}1\_\allowbreak{}a\_\allowbreak{}1\_\allowbreak{}a\_\allowbreak{}06\newline Тип: integer\newline Значення: \textbf{30} \\ \cline{3-5}
 &  & Політика контролю доступу <AC-01\_\allowbreak{}ODP[03] ВИБІРКОВЕ ЗНАЧЕННЯ ПАРАМЕТРА(ів)> містить сферу застосування &  & Параметр: ac\_\allowbreak{}1\_\allowbreak{}a\_\allowbreak{}1\_\allowbreak{}a\_\allowbreak{}07\newline Тип: integer\newline Значення: \textbf{30} \\ \cline{3-5}
 &  & Політика контролю доступу <AC-01\_\allowbreak{}ODP[03] ВИБРАНЕ ЗНАЧЕННЯ ПАРАМЕТРА(ів)> відповідає чинним законам, виконавчим наказам, директивам, положенням, політикам, стандартам та настановам &  & Параметр: ac\_\allowbreak{}1\_\allowbreak{}a\_\allowbreak{}1\_\allowbreak{}b\newline Тип: integer\newline Значення: \textbf{30} \\ \cline{3-5}
 &  & <AC-01\_\allowbreak{}ODP[04] посадова особа> призначається для управління розробкою, документуванням та розповсюдженням політики та процедур контролю доступу &  & Параметр: ac\_\allowbreak{}1\_\allowbreak{}b\newline Тип: integer\newline Значення: \textbf{30} \\ \cline{3-5}
 &  & Переглядається та оновлюється поточна політика контролю доступу <AC-01\_\allowbreak{}ODP[05] частота> &  & Параметр: ac\_\allowbreak{}1\_\allowbreak{}c\_\allowbreak{}01\_\allowbreak{}1\newline Тип: integer\newline Значення: \textbf{30} \\ \cline{3-5}
 &  & Поточну політику контролю доступу переглянуто та оновлено після <AC-01\_\allowbreak{}ODP[06] подій> &  & Параметр: ac\_\allowbreak{}1\_\allowbreak{}c\_\allowbreak{}01\_\allowbreak{}2\newline Тип: integer\newline Значення: \textbf{30} \\ \cline{3-5}
 &  & Переглядаються та оновлюються поточні процедури контролю доступу <AC-01\_\allowbreak{}ODP[07] частота> &  & Параметр: ac\_\allowbreak{}1\_\allowbreak{}c\_\allowbreak{}02\_\allowbreak{}1\newline Тип: integer\newline Значення: \textbf{30} \\ \cline{3-5}
 &  & Поточні процедури контролю доступу переглядаються та оновлюються після <AC-01\_\allowbreak{}ODP[08] подій> &  & Параметр: ac\_\allowbreak{}1\_\allowbreak{}c\_\allowbreak{}02\_\allowbreak{}2\newline Тип: integer\newline Значення: \textbf{30} \\ \cline{3-5}
 &  & Вибрано одне або декілька з наступних ЗНАЧЕНЬ ПАРАМЕТРІВ: \{рівень організації; & AT-1 & Параметр: at\_\allowbreak{}1\_\allowbreak{}odp\_\allowbreak{}03\newline Тип: list\newline Значення: \textbf{рівень організації, рівень місії/бізнес-процесу, рівень системи} \\ 
 &  & рівень місії/бізнес-процесу; &  &  \\ 
 &  & рівень системи\} &  &  \\ \cline{3-5}
 &  & Розроблено та задокументовано політику аудиту та підзвітності & AU-1 & Параметр: au\_\allowbreak{}1\_\allowbreak{}a\_\allowbreak{}01\newline Тип: list\newline Значення: \textbf{default\_\allowbreak{}deny\_\allowbreak{}rule, abac\_\allowbreak{}rule\_\allowbreak{}1} \\ \cline{3-5}
 &  & Політика аудиту та підзвітності доведена до персонал або ролі &  & Параметр: au\_\allowbreak{}1\_\allowbreak{}a\_\allowbreak{}02\newline Тип: list\newline Значення: \textbf{admin, security\_\allowbreak{}officer} \\ \cline{3-5}
 &  & Розроблені та задокументовані процедури аудиту та підзвітності, що сприяють впровадженню політики аудиту та підзвітності, а також відповідні заходи контролю аудиту та підзвітності &  & Параметр: au\_\allowbreak{}1\_\allowbreak{}a\_\allowbreak{}03\newline Тип: list\newline Значення: \textbf{default\_\allowbreak{}deny\_\allowbreak{}rule, abac\_\allowbreak{}rule\_\allowbreak{}1} \\ \cline{3-5}
 &  & Посадова особа призначається для управління політикою та процедурами аудиту та підзвітності AU-01(c)[01][01] переглядається та оновлюється поточна політика аудиту та підзвітності з частота; &  & Параметр: au\_\allowbreak{}1\_\allowbreak{}b\newline Тип: list\newline Значення: \textbf{admin, security\_\allowbreak{}officer} \\ 
 &  & AU-01(c)[01][02] переглядається та оновлюється поточна політика аудиту та підзвітності після подій; &  &  \\ 
 &  & AU-01(c)[02][01] переглядається та оновлюється поточні процедури аудиту та підзвітності з частота; &  &  \\ 
 &  & AU-01(c)[02][02] переглядається та оновлюється поточні процедури аудиту та підзвітності після подій &  &  \\ \cline{3-5}
 &  & Визначено персонал або ролі, до яких має бути доведена політика аудиту та підзвітності &  & Параметр: au\_\allowbreak{}1\_\allowbreak{}odp\_\allowbreak{}01\newline Тип: list\newline Значення: \textbf{admin, security\_\allowbreak{}officer} \\ \cline{3-5}
 &  & Визначено персонал або ролі, на які поширюються процедури аудиту та підзвітності &  & Параметр: au\_\allowbreak{}1\_\allowbreak{}odp\_\allowbreak{}02\newline Тип: list\newline Значення: \textbf{admin, security\_\allowbreak{}officer} \\ \cline{3-5}
 &  & Визначено посадову особу, яка управлятиме політикою та процедурами аудиту та підзвітності &  & Параметр: au\_\allowbreak{}1\_\allowbreak{}odp\_\allowbreak{}04\newline Тип: list\newline Значення: \textbf{admin, security\_\allowbreak{}officer} \\ \cline{3-5}
 &  & Визначено частоту, з якою переглядається та оновлюється поточна політика аудиту та підзвітності &  & Параметр: au\_\allowbreak{}1\_\allowbreak{}odp\_\allowbreak{}05\newline Тип: list\newline Значення: \textbf{default\_\allowbreak{}deny\_\allowbreak{}rule, abac\_\allowbreak{}rule\_\allowbreak{}1} \\ \cline{3-5}
 &  & Визначено події, які потребують перегляду та оновлення поточної політики аудиту та підзвітності &  & Параметр: au\_\allowbreak{}1\_\allowbreak{}odp\_\allowbreak{}06\newline Тип: list\newline Значення: \textbf{default\_\allowbreak{}deny\_\allowbreak{}rule, abac\_\allowbreak{}rule\_\allowbreak{}1} \\ \cline{3-5}
 &  & Визначено частоту, з якою переглядаються та оновлюються поточні процедури аудиту та підзвітності &  & Параметр: au\_\allowbreak{}1\_\allowbreak{}odp\_\allowbreak{}07\newline Тип: integer\newline Значення: \textbf{30} \\ \cline{3-5}
 &  & Визначено події, які потребують перегляду та оновлення поточної процедури аудиту та підзвітності &  & Параметр: au\_\allowbreak{}1\_\allowbreak{}odp\_\allowbreak{}08\newline Тип: list\newline Значення: \textbf{login, logout, failed\_\allowbreak{}attempt} \\ \cline{3-5}
 &  & визначено персонал або ролі, серед яких має бути поширена політика оцінювання, авторизації та моніторингу; & CA-1 & Параметр: ca\_\allowbreak{}1\_\allowbreak{}odp\_\allowbreak{}1\newline Тип: list\newline Значення: \textbf{admin} \\ \cline{3-5}
 &  & визначено персонал або ролі, серед яких мають бути поширені процедури оцінювання, авторизації та моніторингу; &  & Параметр: ca\_\allowbreak{}1\_\allowbreak{}odp\_\allowbreak{}2\newline Тип: list\newline Значення: \textbf{admin} \\ \cline{3-5}
 &  & вибрано одне або декілька з наступних ЗНА ЧЕНЬ ПАРАМЕТРІВ: \{рівень організації; &  & Параметр: ca\_\allowbreak{}1\_\allowbreak{}odp\_\allowbreak{}3\newline Тип: integer\newline Значення: \textbf{30} \\ 
 &  & рівень місії/бізнес -процесу; &  &  \\ 
 &  & рівень системи\}; &  &  \\ \cline{3-5}
 &  & визначено частоту, з якою переглядається та оновлюється поточна політика оцінювання, авторизації та моніторингу; &  & Параметр: ca\_\allowbreak{}1\_\allowbreak{}odp\_\allowbreak{}5\newline Тип: integer\newline Значення: \textbf{30} \\ \cline{3-5}
 &  & визначено частоту, з якою переглядається та оновлюється поточні процедури оцінювання, авторизації та моніторингу; &  & Параметр: ca\_\allowbreak{}1\_\allowbreak{}odp\_\allowbreak{}7\newline Тип: integer\newline Значення: \textbf{30} \\ \cline{3-5}
 &  & розроблені та задокументовані процедури оцінювання, авторизації та моніторингу, що сприяють впровадженню політики оцінювання, авторизації та моніторингу, а також пов'язані з ними засоби контролю оцінювання, авторизації та моніторингу; &  & Параметр: ca\_\allowbreak{}1\_\allowbreak{}a\_\allowbreak{}3\newline Тип: list\newline Значення: \textbf{admin} \\ \cline{3-5}
 &  & процедури оцінювання, авторизації та  моніторингу поширюються серед <CA-01\_\allowbreak{}ODP[02] персоналу або ролей>; &  & Параметр: ca\_\allowbreak{}1\_\allowbreak{}a\_\allowbreak{}4\newline Тип: list\newline Значення: \textbf{admin} \\ \cline{3-5}
 &  & політика <CA-01\_\allowbreak{}ODP[03] ВИБРАНЕ ЗНАЧЕННЯ ПАРАМЕТРА(ів)> оцінювання, авторизації та моніторингу містить ролі; &  & Параметр: ca\_\allowbreak{}1\_\allowbreak{}a\_\allowbreak{}1\_\allowbreak{}a\_\allowbreak{}3\newline Тип: list\newline Значення: \textbf{admin} \\ \cline{3-5}
 &  & політика <CA-01\_\allowbreak{}ODP[03] ВИБРАНЕ ЗНАЧЕННЯ ПАРАМЕТРА(ів)> оцінювання, авторизації та моніторингу містить систему контролю відповідності; &  & Параметр: ca\_\allowbreak{}1\_\allowbreak{}a\_\allowbreak{}1\_\allowbreak{}a\_\allowbreak{}7\newline Тип: list\newline Значення: \textbf{admin} \\ \cline{3-5}
 &  & & CM-1 & \\ \cline{3-5}
 &  & & IA-1 & \\ \cline{3-5}
 &  & & IR-1 & \\ \cline{3-5}
 &  & & MA-1 & \\ \cline{3-5}
 &  & & MP-1 & \\ \cline{3-5}
 &  & & PE-1 & \\ \cline{3-5}
 &  & & PL-1 & \\ \cline{3-5}
 &  & & PS-1 & \\ \cline{3-5}
 &  & & RA-1 & \\ \cline{3-5}
 &  & & SA-1 & \\ \cline{3-5}
 &  & Визначено персонал або ролі, до яких має бути доведена політика захисту системи та комунікацій & SC-1 & Параметр: sc\_\allowbreak{}1\_\allowbreak{}odp\_\allowbreak{}01\newline Тип: list\newline Значення: \textbf{admin, security\_\allowbreak{}officer} \\ \cline{3-5}
 &  & Визначено персонал або ролі, на які поширюються процедури захисту системи та комунікацій &  & Параметр: sc\_\allowbreak{}1\_\allowbreak{}odp\_\allowbreak{}02\newline Тип: list\newline Значення: \textbf{admin, security\_\allowbreak{}officer} \\ \cline{3-5}
 &  & Визначено посадову особу, яка керуватиме політикою та процедурами захисту системи та комунікацій &  & Параметр: sc\_\allowbreak{}1\_\allowbreak{}odp\_\allowbreak{}04\newline Тип: list\newline Значення: \textbf{admin, security\_\allowbreak{}officer} \\ \cline{3-5}
 &  & Визначена періодичність перегляду та оновлення поточної політики захисту системи та комунікацій &  & Параметр: sc\_\allowbreak{}1\_\allowbreak{}odp\_\allowbreak{}05\newline Тип: list\newline Значення: \textbf{default\_\allowbreak{}deny\_\allowbreak{}rule, abac\_\allowbreak{}rule\_\allowbreak{}1} \\ \cline{3-5}
 &  & Є події, які вимагають перегляду та оновлення поточної політики захисту системи та комунікацій &  & Параметр: sc\_\allowbreak{}1\_\allowbreak{}odp\_\allowbreak{}06\newline Тип: list\newline Значення: \textbf{default\_\allowbreak{}deny\_\allowbreak{}rule, abac\_\allowbreak{}rule\_\allowbreak{}1} \\ \cline{3-5}
 &  & Визначена періодичність перегляду та оновлення поточних процедур захисту системи та засобів зв'язку &  & Параметр: sc\_\allowbreak{}1\_\allowbreak{}odp\_\allowbreak{}07\newline Тип: string\newline Значення: \textbf{щорічно} \\ \cline{3-5}
 &  & Визначено персонал або ролі, до яких має бути доведена політика цілісності системи та інформації & SI-1 & Параметр: si\_\allowbreak{}1\_\allowbreak{}odp\_\allowbreak{}01\newline Тип: list\newline Значення: \textbf{admin, security\_\allowbreak{}officer} \\ \cline{3-5}
 &  & Визначено персонал або ролі, на які поширюються процедури цілісності системи та інформації  &  & Параметр: si\_\allowbreak{}1\_\allowbreak{}odp\_\allowbreak{}02\newline Тип: list\newline Значення: \textbf{admin, security\_\allowbreak{}officer} \\ \cline{3-5}
 &  & Визначено посадову особу, відповідальну за управління системою та політикою і процедурами цілісності інформації &  & Параметр: si\_\allowbreak{}1\_\allowbreak{}odp\_\allowbreak{}04\newline Тип: list\newline Значення: \textbf{admin, security\_\allowbreak{}officer} \\ \cline{3-5}
 &  & Визначено періодичність перегляду та оновлення поточної політики цілісності системи та інформації &  & Параметр: si\_\allowbreak{}1\_\allowbreak{}odp\_\allowbreak{}05\newline Тип: list\newline Значення: \textbf{default\_\allowbreak{}deny\_\allowbreak{}rule, abac\_\allowbreak{}rule\_\allowbreak{}1} \\ \cline{3-5}
 &  & Є події, які вимагають перегляду та оновлення поточної політики цілісності системи та інформації &  & Параметр: si\_\allowbreak{}1\_\allowbreak{}odp\_\allowbreak{}06\newline Тип: list\newline Значення: \textbf{default\_\allowbreak{}deny\_\allowbreak{}rule, abac\_\allowbreak{}rule\_\allowbreak{}1} \\ \cline{3-5}
 &  & Визначено частоту, з якою переглядаються та оновлюються поточні цілісності системи та інформації &  & Параметр: si\_\allowbreak{}1\_\allowbreak{}odp\_\allowbreak{}07\newline Тип: integer\newline Значення: \textbf{30} \\ \cline{3-5}
 &  & & SR-1 & \\ \hline
\multicolumn{5}{|>{\raggedright\arraybackslash}p{16cm}|}{\textbf{ПЛАНУВАННЯ (PL)}} \\* \noalign{\hrule height 0.4pt}
90 & ПЛАНИ ЗАХИСТУ ІНФОРМАЦІЇ ТА ПЕРСОНАЛЬНИХ ДАНИХ (PL-2) & & PL-2 & \\ \hline
91 & ПРАВИЛА ПОВЕДІНКИ (PL-4) & & PL-4 & \\ \hline
\multicolumn{5}{|>{\raggedright\arraybackslash}p{16cm}|}{\textbf{ПРИДБАННЯ СИСТЕМ ТА ПОСЛУГ (SA)}} \\* \noalign{\hrule height 0.4pt}
92 & БЕЗПЕКА ТА ПРИВАТНІСТЬ ПРИНЦИПІВ ІНЖИНІРИНГУ (SA-8) & & SA-8 & \\ \hline
93 & КОМПОНЕНТИ СИСТЕМИ, ЩО НЕ ПІДТРИМУЮТЬСЯ (SA-22) & & SA-22 & \\ \hline
94 & ЗОВНІШНІ ПОСЛУГИ ДЛЯ СИСТЕМИ (SA-9) & & SA-9 & \\ \hline
\multicolumn{5}{|>{\raggedright\arraybackslash}p{16cm}|}{\textbf{УПРАВЛІННЯ РИЗИКАМИ В ЛАНЦЮГУ ПОСТАЧАННЯ (SR)}} \\* \noalign{\hrule height 0.4pt}
95 & ПЛАН УПРАВЛІННЯ РИЗИКАМИ ЛАНЦЮГА ПОСТАЧАННЯ (SR-2) & & SR-2 & \\ \hline
96 & СТРАТЕГІЇ ПРИДБАННЯ, ІНСТРУМЕНТИ І МЕТОДИ (SR-5) & & SR-5 & \\ \hline
97 & КОНТРОЛЬ ЛАНЦЮГА ПОСТАЧАННЯ І ПРОЦЕСІВ (SR-3) & & SR-3 & \\ \hline


\end{longtable}

\end{document}
